\documentclass[9pt,french]{article}
\usepackage{verbatim}
\usepackage{amsmath,amsfonts,amssymb}
\usepackage{pgfplots}
\usepackage[utf8]{inputenc}
\usepackage[french]{babel}
\usepackage{pgf,tikz,tkz-tab}
\usepackage{mathrsfs}
\usetikzlibrary{arrows}
\usepackage{pstricks}
\pgfplotsset{compat=newest}
\usepackage{geometry}
\usepackage{eurosym}
\geometry{hmargin=0.5cm,vmargin=1cm}
\pagestyle{empty}
\begin{document}
\underline{Devoir surveillé II}\\
\vspace{1mm}



\underline{La feuille est à rendre avec votre copie. Lors du tracé de vecteur, il est indispensable de marqué le nom du vecteur représenté.}

Nom:\\
Prénom:

\vspace{2mm}




\underline{Exercice 1}

\textit{\underline{Définition}}
\vspace{2mm}

1) Donner les trois critères qui déterminent entièrement un vecteur.
\vspace{2mm}



\textit{\underline{Addition et soustraction de vecteurs}}
\vspace{2mm}


\definecolor{yqyqyq}{rgb}{0.5019607843137255,0.5019607843137255,0.5019607843137255}
\begin{center}

\begin{tikzpicture}[scale=0.6,line cap=round,line join=round,>=triangle 45,x=1.0cm,y=1.0cm]
\draw [color=yqyqyq,, xstep=1.0cm,ystep=1.0cm] (0.6,-4.7) grid (15.92,6.92);
\clip(0.6,-4.7) rectangle (15.92,6.92);
\draw [->,line width=2pt] (3.,0.) -- (5.,2.);
\draw [->,line width=2pt] (10.,3.) -- (6.,3.);
\draw [->,line width=2pt] (6.,1.) -- (8.,1.);
\begin{scriptsize}
\draw[color=black] (3.8,1.24) node {$u$};
\draw[color=black] (7.9,3.3) node {$v$};
\draw[color=black] (6.84,1.3) node {$w$};
\draw[color=black] (13.14,-1.64) node {$A$};
\draw [fill=black] (13.,-2.) circle (2.5pt);
\end{scriptsize}
\end{tikzpicture}

\end{center}


2) $\overrightarrow{w}$ à la même direction que ..... .

3) On pose $\overrightarrow{z}=\overrightarrow{u}+\overrightarrow{v}$
, tracer $\overrightarrow{z}$ en laissant les traits de constructions.

4) On pose $\overrightarrow{a}=\overrightarrow{w}-\overrightarrow{u}$
, tracer $\overrightarrow{a}$ en laissant les traits de constructions. 

5) Est-ce que $\overrightarrow{z}$ et $\overrightarrow{a}$ ont la même direction? 

6) Si oui,  ont-ils même sens?

7)Est-ce que $\overrightarrow{z}$ et $\overrightarrow{a}$ ont même longueur?

8) Tracer  $2\overrightarrow{z}$ et $- \overrightarrow{w}$ et $2\overrightarrow{z}+\overrightarrow{u}$.

9) Déterminer le point $M$ tel que $\overrightarrow{AM}=\overrightarrow{z}$

\vspace{2mm}



\textit{\underline{Relation de Chasles}}

\vspace{2mm}



On rappel la relation de Chasles: 

$$ \overrightarrow{AB} + \overrightarrow{BC} = \overrightarrow{AC} $$
\definecolor{xfqqff}{rgb}{0.4980392156862745,0.,1.}
\definecolor{ffqqqq}{rgb}{1.,0.,0.}
\definecolor{qqqqff}{rgb}{0.,0.,1.}

\begin{center}

\begin{tikzpicture}[scale=0.4,line cap=round,line join=round,>=triangle 45,x=1.0cm,y=1.0cm]
\clip(0.58,-4.72) rectangle (15.9,6.9);
\draw [->,color=ffqqqq] (1.54,-0.86) -- (15.3,0.02);
\draw [->,color=qqqqff] (15.3,0.02) -- (5.54,3.52);
\draw [->,color=xfqqff] (1.54,-0.86) -- (5.54,3.52);
\begin{scriptsize}
\draw [fill=qqqqff] (1.54,-0.86) circle (2.5pt);
\draw[color=qqqqff] (1.3,-0.5) node {$A$};
\draw [fill=qqqqff] (15.3,0.02) circle (2.5pt);
\draw[color=qqqqff] (15.44,0.38) node {$B$};
\draw [fill=qqqqff] (5.54,3.52) circle (2.5pt);
\draw[color=qqqqff] (5.8,3.88) node {$C$};
\end{scriptsize}
\end{tikzpicture}

\end{center}
10) Réduire les expressions vectorielles suivantes à l'aide de la relation de Chasles:




\begin{itemize}
\item[a)]   $\overrightarrow{AB} + \overrightarrow{MA} + \overrightarrow{CM}$
\item[b)] $\overrightarrow{AB} + \overrightarrow{MA} - \overrightarrow{MC}$
\item[c)] $\overrightarrow{CA} + \overrightarrow{AB} + \overrightarrow{BC}$
\end{itemize}


\underline{Exercice 2}\\

Soient $ABCD$ un parallélogramme de centre $O$ et $\mathcal{C}$ le cercle de centre $O$ de rayon $OB$.\\



\underline{L'objectif est de montrer que $EBFD$ est un rectangle}\\

\definecolor{dbwrru}{rgb}{0.8588235294117647,0.3803921568627451,0.0784313725490196}
\definecolor{qqwuqq}{rgb}{0,0.39215686274509803,0}
\definecolor{sexdts}{rgb}{0.1803921568627451,0.49019607843137253,0.19607843137254902}
\definecolor{dtsfsf}{rgb}{0.8274509803921568,0.1843137254901961,0.1843137254901961}
\definecolor{wrwrwr}{rgb}{0.3803921568627451,0.3803921568627451,0.3803921568627451}
\definecolor{rvwvcq}{rgb}{0.08235294117647059,0.396078431372549,0.7529411764705882}
\definecolor{cqcqcq}{rgb}{0.7529411764705882,0.7529411764705882,0.7529411764705882}
\begin{center}
\begin{tikzpicture}[scale=0.7,line cap=round,line join=round,>=triangle 45,x=1cm,y=1cm]
\draw [color=cqcqcq,, xstep=2cm,ystep=2cm] (-7.106889489791369,-4.935186427211233) grid (17.91185394507301,12.510184962089731);
\clip(-7.106889489791369,-4.935186427211233) rectangle (17.91185394507301,12.510184962089731);
\draw [line width=2pt,color=wrwrwr] (-6,-2)-- (2,-2);
\draw [line width=2pt,color=wrwrwr] (8,8)-- (16,8);
\draw [line width=2pt,color=wrwrwr] (2,-2)-- (8,8);
\draw [line width=2pt,color=wrwrwr] (-6,-2)-- (16,8);
\draw [line width=0.8pt,dash pattern=on 1pt off 1pt on 1pt off 4pt,color=dtsfsf] (5,3) circle (5.8309518948453cm);
\draw [line width=2pt,color=dbwrru] (-0.3083031367642161,0.5871349378344473)-- (2,-2);
\draw [line width=2pt,color=dbwrru] (2,-2)-- (10.308303136764216,5.412865062165553);
\draw [line width=2pt,color=dbwrru] (10.308303136764216,5.412865062165553)-- (8,8);
\draw [line width=2pt,color=dbwrru] (8,8)-- (-0.3083031367642161,0.5871349378344473);
\draw [line width=2pt,color=wrwrwr] (-6,-2)-- (8,8);
\draw [line width=2pt,color=wrwrwr] (2,-2)-- (16,8);
\begin{scriptsize}
\draw [fill=rvwvcq] (-6,-2) circle (2.5pt);
\draw[color=rvwvcq] (-5.824497418178399,-1.5598337104186057) node {$A$};
\draw [fill=rvwvcq] (8,8) circle (2.5pt);
\draw[color=rvwvcq] (8.16083498544985,8.457342094256289) node {$D$};
\draw [fill=rvwvcq] (2,-2) circle (2.5pt);
\draw[color=rvwvcq] (2.2206981254312774,-1.4993435183613901) node {$B$};
\draw [fill=rvwvcq] (16,8) circle (2.5pt);
\draw[color=rvwvcq] (16.169736413825195,8.457342094256289) node {$C$};
\draw [fill=dtsfsf] (5,3) ++(-1.5pt,0 pt) -- ++(1.5pt,1.5pt)--++(1.5pt,-1.5pt)--++(-1.5pt,-1.5pt)--++(-1.5pt,1.5pt);
\draw[color=dtsfsf] (4.90521459411956,3.351969884627297) node {$O$};
\draw[color=dtsfsf] (1.6762863969163366,8.408949940610517) node {$\mathcal{C}$};
\draw [color=sexdts] (-0.3083031367642161,0.5871349378344473)-- ++(-2pt,-2pt) -- ++(4pt,4pt) ++(-4pt,0) -- ++(4pt,-4pt);
\draw[color=sexdts] (-0.13841936480013156,0.9807543559844476) node {$E$};
\draw [color=qqwuqq] (10.308303136764216,5.412865062165553)-- ++(-2pt,-2pt) -- ++(4pt,4pt) ++(-4pt,0) -- ++(4pt,-4pt);
\draw[color=qqwuqq] (10.48365836044693,5.819969720561692) node {$F$};
\end{scriptsize}
\end{tikzpicture}
\end{center}
Compléter le texte à trou suivant:\\

Comme $O$ est le milieu de $[DB]$ alors $\overrightarrow{\cdots}=\overrightarrow{\cdots}$.

Comme $[EF]$ est un ............... du cercle $\mathcal{C}$ alors $O$ est le milieu de $[....]$, donc $\overrightarrow{\cdots}=\overrightarrow{\cdots}$. 

Comme $[EF]$ et $[DB]$ sont de même ................... . \\

Ainsi comme $EBFD$ est un quadrilatère dont ces diagonales $[EF]$ et $[DB]$ se coupent en leurs milieux ...., $EBFD$ est un .............

De plus comme $EF=BD$ alors $EBFD$ est un .................. . 


\end{document}


